\section{Notation and Package Model}

Let an input asset contain source shapes indexed by $s \in S$ and candidate primitive colliders
indexed by $i \in I$. A candidate primitive package is written
\[
  p_i = (k_i, F_i, T_i, \theta_i, \ell_i),
\]
where $k_i$ is the primitive kind, $F_i$ is the source frame, $T_i$ is the primitive transform,
$\theta_i$ stores dimensions and material-side metadata, and $\ell_i$ is the owning link identity
when one is available.

\begin{table}[t]
\centering
\caption{Supplement-only notation table. It defines audit terms used by the derivations and does
not duplicate a main-paper result table.}
\label{tab:supp-notation}
\begin{tabular}{ll}
\toprule
Symbol & Meaning \\
\midrule
$S$ & source collision or visual shapes attached to an asset \\
$I$ & candidate primitive collider indices \\
$p_i$ & one generated primitive collider candidate \\
$B$ & simulated rigid bodies or robot links \\
$a(p_i)$ & body attachment map for a primitive candidate \\
$D$ & diagnostic scene with probes and acceptance predicates \\
\bottomrule
\end{tabular}
\end{table}

For a link-preserving package, the attachment map is constrained by
\[
  a(p_i) = b_{\ell_i}.
\]
The equality is a structural requirement: a generated primitive for one link must not be silently
attached to a neighboring link just because the combined geometry looks plausible in isolation.
